% Esquisse d'un programme (cote 119, copy 1): archive cote 119, site page 33 (PDF page 34), typescript page 31.
% typescript. Draft compared: EsquisseFr.pdf (text layer):1129-1190.
% AI-transcribed by Claude Opus 5.5 (claude-opus-5-5, adaptive thinking, effort high), Message Batches API, batch esq1 (msgbatch_01LNsogFyF3y6KSp9UbtJtBk), 2026-11-12; not yet checked by eye.
\scanpage{119/33}
\mspage{31}

stricte de toutes, celles des ``espaces $\bar{\mathbb{Q}}_r$-algébriques par morceaux''
(où $\bar{\mathbb{Q}}_r = \bar{\mathbb{Q}} \cap \mathbb{R}$), vers celle qui (à tort ou à raison) m'apparaît
comme probablement la plus vaste, savoir celle des ``espaces ana-%
litiques réels par morceaux'' (ou semianalytiques dans la termi-%
nologie de Hironaka). Parmi les théorèmes de fondements envisagés
dans mon programme, il y a un \uline{théorème de comparaison} qui, vague-%
ment parlant, dira qu'\uline{on trouvera essentiellement les mêmes caté-%
gories isotopiques} (ou même $\infty$-isotopiques), quelle que soit la
théorie modérée avec laquelle on travaille\hand{$^{(6)}$}. De façon plus précise,
il s'agit de mettre le doigt sur un système d'axiomes suffisam-%
ment riche, pour impliquer (entre bien autres choses !) que si on
a deux théories modérées \hand{$\mathcal{T}$}, \hand{$\mathcal{T}$}' avec \hand{$\mathcal{T}$} plus fine que \hand{$\mathcal{T}$}' (dans
un sens évident), et si X, Y sont deux espaces T-modérés, qui
définissent donc aussi des espaces T'-modérés correspondants,
l'application canonique
\[ \underline{\mathrm{Isom}}_{\hand{\mathcal{T}}}(X,Y) \longrightarrow \underline{\mathrm{Isom}}_{\hand{\mathcal{T}}'}(X,Y) \]
induit une bijection sur l'ensemble des composantes connexes (ce
qui impliquera que la catégorie isotopique des T-espaces est équi-%
valente à celle des T'-espaces), et même, est une équivalence
d'homotopie (ce qui signifie qu'on a même une équivalence pour
les catégories ``$\infty$-isotopique'', plus fines que les catégories
isotopiques où on ne retient que le $\pi_0$ des espaces d'isomorphis-%
mes). Ici les \uline{Isom} peuvent être définis de façon évidente comme
ensembles semisimpliciaux par exemple, pour pouvoir donner un sens
précis à l'énoncé précédent. Des énoncés analogues devraient être
vrais, en remplaçant les ``espaces'' \uline{Isom} par d'autres espaces d'ap-%
plications, soumises à des conditions géométriques standard, comme
celle d'être des plongements, des immersions, lisses, étales, des
fibrations etc. Egalement, on s'attend à avoir des énoncés analo-%
gues, où X, Y sont remplacés par des systèmes d'espaces modérés,
tels ceux qui interviennent dans une théorie de dévissage des
structures stratifiées - de telle sorte que dans un sens techni-%
que précis, cette théorie de dévissage sera, elle aussi, essen-%

\margin{33}
